(1) 棒两端的动生电动势大小为 \(\mathcal E=B\ell v\)，电流大小 \(I=B\ell v/R\)。根据楞次定律，电流产生的电磁力必须阻碍棒的前进，所以
\[
F_x=-I\ell B=-\frac{B^2\ell^2}{R}v,\qquad
m\dot v=-\frac{B^2\ell^2}{R}v.
\]

(2) 定义 \(\tau=mR/(B^2\ell^2)\)，满足一阶指数衰减方程，故
\[
v(t)=v_0e^{-t/\tau},\quad
x(t)-x_0=v_0\tau(1-e^{-t/\tau}),\quad
I(t)=\frac{B\ell v_0}{R}e^{-t/\tau}.
\]
该结果同时给定棒的机械轨迹和瞬时电流。

(3) 电阻的瞬时热功率 \(P_R=I^2R=B^2\ell^2v_0^2e^{-2t/\tau}/R\)。因此
\[
Q_R=\int_0^\infty P_Rdt
=\frac{B^2\ell^2v_0^2}{R}\frac{\tau}{2}
=\frac12mv_0^2.
\]
或者将动力学方程乘以 v 得 \(\frac d{dt}(mv^2/2)=-I^2R\)，积分后有相同结果。全部初始动能最终变成电阻焦耳热。

(4) 棒在理想无限长导轨上最终位移是 \(x_\infty-x_0=v_0\tau\)。当 \(R\to \infty \) 或 \(B\to 0\) 时 \(\tau \to \infty \)，对任何固定有限 t，\(v\to v_0\)、\(I\to0\)，位移趋于 \(v_{0}t\)，而不是已经停止。必须区分先令 \(t\to \infty \) 与先令阻尼消失的极限；磁场通过洛伦兹力提供机械阻力，而电路电场对载流子做功并将能量转化为热，符合完整能量预算。
